\section{结论}

本轮特征分析已经完成两个数据集、三个任务和两个特征集方向的收敛。

最终决策：

\begin{center}
\config{manual_basic} 是当前 downstream baseline 的主特征集。
\end{center}

任务级结论：

\begin{itemize}
  \item RUL：优先使用振幅/能量类时域特征。
  \item Health State：优先使用幅值类特征，XJTU-SY 需考虑 cross-condition 稳定性，PHM2012 更偏 horizontal 通道。
  \item Early Fault：PHM2012 仍由幅值类特征主导；XJTU-SY 需要明确工况敏感 caveat。
\end{itemize}

\feature{mag__time__rms} 保留为 label-source reference，不作为独立发现证据。

下一步应进入 baseline planning。后续实验应基于 \config{manual_basic} 和最终推荐表：

\begin{center}
\path{../summary/recommended_features.csv}
\end{center}

然后分别设计 RUL 回归、Health State 分类和 Early Fault Detection 的可复现实验。
